% =============================================================================
%  本书概述（原主文件里的封底文案 + 作者简介 + 本书相关）
%  · 引文用 \myQuote（灰底细框），与原书 Quote Box 一致
% =============================================================================

完整的介绍 C++ 移动语义。

C++11 添加的移动语义已经成为现代 C++ 的标志，也使语言变得复杂，即使经验丰富的开发者仍需要仔细处理。所以，一些编程书籍甚至不推荐对非常简单的类使用移动语义。因此，详细地解释 C++ 移动语义就变得刻不容缓。

本书会从基本原理开始来介绍移动语义，并解释移动语义的细节，使每个开发者都可以正确地使用移动语义。

你将学习到：
\begin{itemize}
  \item 移动语义的起因和术语
  \item 如何隐式地获益于移动语义
  \item 如何显式地获益于移动语义
  \item 会遇到的问题，以及如何处理它们
  \item 所有的结果都取决于你的编程风格
\end{itemize}

重点在于所描述的特性，需要在实践中进行应用。示例和背景信息，有助于理解和改进简单类，甚至泛型库和框架的代码。

\begin{myQuote}
“我以为我理解了移动的语义，其实我真的不懂！我从你的书中学到了很多东西。”\\
\hfill —— Jonathan Boccara
\end{myQuote}

\begin{myQuote}
“这是我非常需要的书。”\\
\hfill —— Rob Bernstein
\end{myQuote}

\begin{myQuote}
“有时候我觉得我对纠缠和量子隐形传态的理解，要比我对一些奇怪的 C++ 移动语义的理解要好。套用 Feynman 的话：如果你认为你理解了 C++ 的移动语义，那你就不理解 C++ 的移动语义。赶快阅读这本书吧。”\\
\hfill —— Victor Ciura
\end{myQuote}

\mySmallsection{作者简介}

Nicolai Josuttis（\texthref{http://www.josuttis.com}）在编程界很有名，其发言和著作都很有权威，还是世界范围内畅销书的（共同）作者：
\begin{itemize}
  \item 《The C++ Standard Library》
  \item 《C++ Templates》
  \item 《C++ Move Semantics》
  \item 《C++17》
  \item 《SOA in Practice》
\end{itemize}

同时也是一位富有创新精神的演讲者，曾在各种会议和活动中发言。还是独立的讲师，并且在 C++ 标准化方面有 20 多年的经验。

\mySmallsection{本书相关}

\begin{itemize}
  \item github 翻译地址：\href{https://github.com/xiaoweiChen/CPP-Move-Semantics}{https://github.com/xiaoweiChen/CPP-Move-Semantics}
\end{itemize}
